% !TEX root = web.tex
% Lean blueprint fragment for the Moise / PL-complex route.
% Intended location: blueprint/src/content_moise_pl.tex
%
% This is a content.tex-style fragment: include it from content.tex with
%   \input{content_moise_pl}
% It assumes the usual leanblueprint macros are available: \lean, \uses,
% \leanok, \mathlibok, \notready, and standard theorem environments.
%
% This file is written as an alternative complete route.  The Gallier--Xu and
% quotient-space tail intentionally mirrors the Mohar--Thomassen route so the two
% files can be evaluated independently.  In a combined blueprint, extract the
% shared tail into one common file to avoid duplicate labels and declarations.

\chapter{Route B: Moise via PL complexes}

\section{Global interface for the classification project}

\begin{definition}[Lean Eval compact surface with boundary]
\label{moise:def:eval_surface}
\lean{LeanEval.Topology.ClassificationOfSurfaces.EvalSurface}
An \emph{Eval surface} is a type $S$ equipped with a topology such that $S$ is
Hausdorff, connected, compact, charted over the Euclidean half-space
$\mathbb H^2$, and is a $C^0$ manifold with corners for
\texttt{modelWithCornersEuclideanHalfSpace 2}.  This definition packages exactly
the hypotheses of the Lean Eval problem.
\end{definition}

\begin{definition}[Finite surface cell complex]
\label{moise:def:surface_cell_complex}
\lean{LeanEval.Topology.ClassificationOfSurfaces.SurfaceCellComplex}
A \emph{finite surface cell complex} is a finite combinatorial presentation of a
compact bordered surface by polygonal faces and oriented edge identifications.
It contains finite types of faces, oriented darts, edge-pairing data, cyclic
boundary words for faces, and incidence conditions ensuring that each vertex
link is either a circle or an interval.
\end{definition}

\begin{definition}[Realization of a surface cell complex]
\label{moise:def:cell_realization}
\lean{LeanEval.Topology.ClassificationOfSurfaces.SurfaceCellComplex.Realization}

\uses{moise:def:surface_cell_complex}
For a finite surface cell complex $K$, the realization $|K|$ is the quotient of
the disjoint union of its standard polygonal faces by the equivalence relation
identifying paired oriented boundary arcs.  It carries the quotient topology.
\end{definition}

\begin{theorem}[Topological bridge]
\label{moise:thm:compact_surface_cellulable}
\lean{LeanEval.Topology.ClassificationOfSurfaces.compact_surface_homeomorphic_to_cell_complex}

\uses{moise:def:eval_surface,moise:def:surface_cell_complex,moise:def:cell_realization,moise:thm:compact_eval_surface_finitely_triangulable,moise:thm:finite_triangulation_to_cell_complex}
Every Eval surface $S$ is homeomorphic to the realization of some finite surface
cell complex: there exists a finite surface cell complex $K$ and a homeomorphism
$S\cong |K|$.
\end{theorem}

\begin{proof}
\notready

Use the Moise/Rado theorem, strengthened to compact manifolds with boundary, to
obtain a finite triangulation of $S$.  Convert that triangulation to a finite
surface cell complex and compose the homeomorphisms.
\end{proof}

\section{PL-complex foundations}

\begin{definition}[Euclidean finite complex and support]
\label{moise:def:euclidean_complex}
\lean{LeanEval.Topology.ClassificationOfSurfaces.EuclideanComplex, LeanEval.Topology.ClassificationOfSurfaces.EuclideanComplex.support}
\leanok
A Euclidean finite complex is a finite simplicial complex made of rectilinear
simplexes in some Euclidean space.  Its support $|K|$ is the union of its
simplexes with the subspace topology.
\end{definition}

\begin{definition}[Subdivision]
\label{moise:def:subdivision}
\lean{LeanEval.Topology.ClassificationOfSurfaces.EuclideanComplex.Subdivision}
\leanok

\uses{moise:def:euclidean_complex}
A subdivision $K'<K$ is a finite complex with homeomorphic support and a carrier
map sending each simplex of $K'$ to a simplex of $K$.  The formal API records
dimension monotonicity, face compatibility, and coarse coverage data; the
geometric rectilinear containment is isolated as refinement data.
\end{definition}

\begin{definition}[PL map and PL homeomorphism]
\label{moise:def:pl_map_homeomorphism}
\lean{LeanEval.Topology.ClassificationOfSurfaces.PLMap, LeanEval.Topology.ClassificationOfSurfaces.PLHomeomorph}
\leanok

\uses{moise:def:euclidean_complex,moise:def:subdivision}
A map $f:|K|\to |L|$ is piecewise linear if, after subdividing $K$, its
restriction to each simplex is affine into a simplex of $L$.  A PL homeomorphism
is a homeomorphism whose map and inverse are PL after suitable subdivisions.  The
formal API now includes identity, composition, inverse, and coherence fields
relating the PL maps to the underlying homeomorphism.
\end{definition}

\begin{lemma}[PL property is invariant under subdivision]
\label{moise:lem:pl_invariant_under_subdivision}
\lean{LeanEval.Topology.ClassificationOfSurfaces.pl_iff_pl_after_subdivision}

\uses{moise:def:pl_map_homeomorphism,moise:def:subdivision}
A homeomorphism is PL relative to $K$ and $L$ if and only if it is PL relative to
any subdivision of the domain or target complex.
\end{lemma}

\begin{proof}
\notready

Pass to a common subdivision of the complexes witnessing piecewise linearity.
For target subdivisions, pull back a common subdivision of the image complex.
\end{proof}

\begin{definition}[Combinatorial two-manifold with boundary]
\label{moise:def:combinatorial_surface}
\lean{LeanEval.Topology.ClassificationOfSurfaces.CombinatorialTwoManifoldWithBoundary}
\leanok

\uses{moise:def:euclidean_complex}
A two-dimensional complex $K$ is a combinatorial two-manifold with boundary if
the link of every vertex is either a combinatorial circle or a combinatorial
interval.  The formal definition now exposes named finite link predicates,
boundary and interior vertex predicates, and one- and two-dimensional simplex
filters.  Its support is then expected to be a topological surface with
boundary.
\end{definition}

\begin{definition}[Combinatorial two-cell]
\label{moise:def:combinatorial_two_cell}
\lean{LeanEval.Topology.ClassificationOfSurfaces.CombinatorialTwoCell}
\leanok

\uses{moise:def:euclidean_complex,moise:def:combinatorial_surface}
A combinatorial two-cell is a finite two-dimensional complex with a distinguished
boundary subcomplex, a boundary inclusion map, and PL homeomorphism data to a
closed triangle model respecting the boundary subcomplex.
\end{definition}

\begin{lemma}[Polygonal disk is a combinatorial two-cell]
\label{moise:lem:polygonal_disk_combinatorial_cell}
\lean{LeanEval.Topology.ClassificationOfSurfaces.polygonal_disk_is_combinatorial_two_cell}

\uses{moise:def:combinatorial_two_cell,moise:def:pl_map_homeomorphism}
If $J$ is a polygonal simple closed curve in $\mathbb R^2$ and $I$ is its
interior, then $\overline I$ has a triangulation making it a combinatorial
two-cell.
\end{lemma}

\begin{proof}
\notready
Triangulate the polygonal disk and prove it is PL-homeomorphic to a standard
closed triangle by repeatedly deleting free triangles or by an explicit fan
triangulation after choosing an interior point.
\end{proof}

\begin{theorem}[PL Schoenflies for combinatorial two-cells]
\label{moise:thm:pl_schoenflies_combinatorial_cell}
\lean{LeanEval.Topology.ClassificationOfSurfaces.pl_schoenflies_combinatorial_two_cell}

\uses{moise:def:combinatorial_two_cell,moise:def:pl_map_homeomorphism}
Every PL homeomorphism between the boundaries of two combinatorial two-cells
extends to a PL homeomorphism between the two cells whose restriction to the
distinguished boundary subcomplexes agrees with the prescribed boundary map.
\end{theorem}

\begin{proof}

Choose PL homeomorphisms from the standard 2-simplex onto the two cells.  Reduce
to a PL homeomorphism of the boundary of the standard simplex and extend it over
the simplex, then conjugate back.
\end{proof}

\section{PL approximation in dimension two}

\begin{definition}[Strongly positive function]
\label{moise:def:strongly_positive}
\lean{LeanEval.Topology.ClassificationOfSurfaces.StronglyPositive}
\leanok
A function $\varphi:X\to\mathbb R$ is strongly positive if $\varphi(x)>0$ for all
$x$ and it has the lower-separation properties needed to make local approximation
estimates uniform on compact simplexes.  In a metric formalization this may be
implemented as a positive continuous function or as Moise's strongly positive
notion.
\end{definition}

\begin{definition}[$\varphi$-approximation]
\label{moise:def:phi_approximation}
\lean{LeanEval.Topology.ClassificationOfSurfaces.PhiApproximation}
\leanok

\uses{moise:def:strongly_positive}
For maps $f,g:X\to Y$ into a metric space, $f$ is a $\varphi$-approximation of
$g$ if $d(f(x),g(x))<\varphi(x)$ for every $x$.
\end{definition}

\begin{lemma}[Approximation calculus]
\label{moise:lem:approximation_calculus}
\lean{LeanEval.Topology.ClassificationOfSurfaces.StronglyPositive.mono, LeanEval.Topology.ClassificationOfSurfaces.StronglyPositive.add, LeanEval.Topology.ClassificationOfSurfaces.StronglyPositive.min, LeanEval.Topology.ClassificationOfSurfaces.PhiApproximation.mono, LeanEval.Topology.ClassificationOfSurfaces.PhiApproximation.symm, LeanEval.Topology.ClassificationOfSurfaces.PhiApproximation.trans}
\leanok

\uses{moise:def:strongly_positive,moise:def:phi_approximation}
Strongly positive tolerances are closed under monotonicity, addition, minimum,
positive scalar multiplication, pullback, and restriction.  Pointwise
$\varphi$-approximations are closed under monotonicity, symmetry, triangle
composition, precomposition, restriction, and suitable postcomposition.
\end{lemma}

\begin{proof}
\leanok
These are direct consequences of positivity and the metric triangle inequality.
\end{proof}

\begin{lemma}[PL approximation on the one-skeleton]
\label{moise:lem:pl_approximation_one_skeleton}
\lean{LeanEval.Topology.ClassificationOfSurfaces.pl_approximation_one_skeleton}
\leanok

\uses{moise:def:combinatorial_surface,moise:def:phi_approximation,moise:def:pl_map_homeomorphism,moise:lem:one_skeleton_hard_boundaries}
Let $K$ be a combinatorial two-manifold with boundary, $h:|K|\to M\subseteq
\mathbb R^2$ a homeomorphism, and $\varphi\gg 0$.  After a sufficiently fine
subdivision, $h$ restricted to the one-skeleton has a PL embedding which is a
$\varphi$-approximation and preserves the images of vertices.
\end{lemma}

\begin{proof}
\leanok
Approximate each edge image by a broken line inside a sufficiently small
neighborhood.  Choose the subdivision fine enough so that nonincident edge images
remain separated and incident edges meet only at their common vertices.
\end{proof}

\begin{lemma}[Hard one-skeleton approximation boundaries]
\label{moise:lem:one_skeleton_hard_boundaries}
\lean{LeanEval.Topology.ClassificationOfSurfaces.finite_arc_polygonal_approximation, LeanEval.Topology.ClassificationOfSurfaces.endpoint_preservation_for_polygonal_approximation, LeanEval.Topology.ClassificationOfSurfaces.finite_edge_separation_control, LeanEval.Topology.ClassificationOfSurfaces.no_crossing_perturbation}

\uses{moise:def:combinatorial_surface,moise:def:phi_approximation}
The proof of one-skeleton approximation is split into named theorem boundaries:
finite polygonal approximation of edge images, endpoint preservation, finite
edge separation control, and no-crossing perturbation of the approximated edge
arcs.
\end{lemma}

\begin{proof}
\notready
These are the remaining planar topology and finite-family perturbation
ingredients behind the formal one-skeleton approximation wrapper.
\end{proof}

\begin{definition}[Cellwise extension and gluing interfaces]
\label{moise:def:cellwise_extension_gluing}
\lean{LeanEval.Topology.ClassificationOfSurfaces.CellwiseExtension, LeanEval.Topology.ClassificationOfSurfaces.ExtensionsAgreeOnSharedBoundary, LeanEval.Topology.ClassificationOfSurfaces.GlobalPLHomeomorphFromCellwise, LeanEval.Topology.ClassificationOfSurfaces.GlobalPLSurfaceApproximation}
\leanok

\uses{moise:def:combinatorial_surface,moise:def:phi_approximation,moise:lem:pl_approximation_one_skeleton}
A cellwise extension records a one-skeleton PL approximation, a global candidate
map on the support, its $\varphi$-closeness to the original homeomorphism,
finite two-cell extension data, and finite agreement conditions on shared
boundaries.  The global gluing predicate packages the claim that compatible
cellwise data assemble into a global PL surface approximation.
\end{definition}

\begin{lemma}[Hard cellwise extension and gluing boundaries]
\label{moise:lem:cellwise_extension_gluing_boundaries}
\lean{LeanEval.Topology.ClassificationOfSurfaces.cellwise_extension_by_pl_schoenflies, LeanEval.Topology.ClassificationOfSurfaces.global_pl_homeomorph_from_cellwise}

\uses{moise:def:cellwise_extension_gluing,moise:thm:pl_schoenflies_combinatorial_cell}
The plane approximation proof is split into two theorem boundaries after the
one-skeleton is controlled: extend each two-cell boundary map by PL Schoenflies,
then glue the compatible finite family of cellwise extensions into a global PL
homeomorphic approximation.
\end{lemma}

\begin{proof}
\notready
The first boundary applies the combinatorial PL Schoenflies theorem face by
face.  The second boundary proves that the finite agreement data on shared
boundaries are sufficient for a global continuous PL homeomorphism and that the
chosen separation bounds prevent overlaps of distinct cell interiors.
\end{proof}

\begin{theorem}[Moise PL approximation theorem in the plane]
\label{moise:thm:pl_approximation_plane}
\lean{LeanEval.Topology.ClassificationOfSurfaces.pl_approximation_plane_combinatorial_surface}
\leanok

\uses{moise:def:combinatorial_surface,moise:def:phi_approximation,moise:lem:pl_approximation_one_skeleton,moise:lem:cellwise_extension_gluing_boundaries}
Let $K$ be a combinatorial two-manifold with boundary, let $h:|K|\to M\subseteq
\mathbb R^2$ be a homeomorphism, and let $\varphi\gg 0$.  There exists a PL
homeomorphism $f:|K|\to\mathbb R^2$ which is a $\varphi$-approximation of $h$.
\end{theorem}

\begin{proof}
\leanok

Approximate $h$ on the one-skeleton by a PL embedding.  Apply the cellwise PL
Schoenflies extension boundary to extend across every two-cell, then apply the
global gluing boundary to assemble the compatible cellwise extensions.
\end{proof}

\begin{theorem}[PL approximation between combinatorial surfaces]
\label{moise:thm:pl_approximation_between_surfaces}
\lean{LeanEval.Topology.ClassificationOfSurfaces.pl_approximation_between_combinatorial_surfaces}

\uses{moise:thm:pl_approximation_plane,moise:def:combinatorial_surface}
If $K_1$ is a combinatorial two-manifold with boundary, $K_2$ is a combinatorial
two-manifold, $h:|K_1|\to |K_2|$ is a homeomorphism, and $\varphi\gg0$, then $h$
has a PL homeomorphic $\varphi$-approximation into $|K_2|$.
\end{theorem}

\begin{proof}
\notready

Work inside star neighborhoods in $K_2$ that are PL-homeomorphic to convex
plane cells.  Apply the plane approximation theorem locally after subdividing
$K_1$ so that each small star maps into one such chart, then glue the local PL
extensions.
\end{proof}

\section{Open subcomplexes and PL complexes in a manifold}

\begin{definition}[PL complex in a topological space]
\label{moise:def:pl_complex_in_space}
\lean{LeanEval.Topology.ClassificationOfSurfaces.PLComplexInSpace}
\leanok

\uses{moise:def:euclidean_complex,moise:def:pl_map_homeomorphism}
A PL complex in a topological space $M$ is a locally finite family of subsets of
$M$ equipped with charts from Euclidean simplexes such that the induced abstract
complex is PL-compatible on intersections.  The formal API stores a mathlib
topological embedding of the complex support into the ambient space, exposes the
image support, point and set coverage predicates, extension data, and
compatibility on overlaps.
\end{definition}

\begin{definition}[Interior subcomplex of a PL surface]
\label{moise:def:pl_complex_interior_subcomplex}
\lean{LeanEval.Topology.ClassificationOfSurfaces.PLComplexInSpace.interiorSubcomplex}
\leanok

\uses{moise:def:pl_complex_in_space,moise:def:combinatorial_surface}
For a PL complex $\mathcal X$ whose support is a two-manifold with boundary,
$\mathcal X^\circ$ is the set of simplexes whose image sets do not meet the
boundary of the support.  Until boundary strata are available as geometric
subcomplexes, the formal carrier is the whole support with an explicit openness
lemma.
\end{definition}

\begin{theorem}[Open subset of a finite complex is a complex]
\label{moise:thm:open_subset_complex}
\lean{LeanEval.Topology.ClassificationOfSurfaces.open_subset_of_finite_complex_is_complex}

\uses{moise:def:euclidean_complex,moise:def:subdivision,moise:def:pl_complex_in_space}
If $K$ is a finite complex and $U\subseteq |K|$ is open, then there is a complex
$K_U$ with $|K_U|=U$ such that every simplex of $K_U$ is a rectilinear subsimplex
of a simplex of $K$.  The formal statement returns open-subset complex data: a
new finite complex, a homeomorphism from its support to $U$, an embedding back
into the original support, and an ambient compatibility predicate.
\end{theorem}

\begin{proof}
\notready
Choose a sequence of subdivisions with mesh tending to zero.  Recursively select
subcomplexes lying in $U$ whose supports exhaust $U$, using compact separation
from $|K|\setminus U$.  Subdivide boundary layers compatibly so the union is a
locally finite complex.
\end{proof}

\begin{lemma}[Locally finite compact support implies finite]
\label{moise:lem:compact_locally_finite_complex_finite}
\lean{LeanEval.Topology.ClassificationOfSurfaces.locallyFiniteComplex_finite_of_compact_support}
\leanok

\uses{moise:def:pl_complex_in_space}
A locally finite PL complex with compact support has only finitely many simplexes.
\end{lemma}

\begin{proof}
\leanok

Each point of the compact support has a neighborhood meeting only finitely many
simplexes.  Extract a finite subcover and take the union of the corresponding
finite sets of simplexes.
\end{proof}

\section{Moise--Rado triangulation}

\begin{definition}[Moise topological two-manifold]
\label{moise:def:moise_two_manifold}
\lean{LeanEval.Topology.ClassificationOfSurfaces.MoiseTwoManifold}
\leanok
A Moise two-manifold is a separable metric space locally homeomorphic to the open
unit disk.  The bordered version is locally homeomorphic to either the open disk
or the closed half-disk.
\end{definition}

\begin{lemma}[Chart pair exhaustion]
\label{moise:lem:chart_pair_exhaustion}
\lean{LeanEval.Topology.ClassificationOfSurfaces.chart_pair_exhaustion}

\uses{moise:def:moise_two_manifold}
Every Moise two-manifold admits a countable sequence of chart pairs
$(N_i,N'_i)$ such that $N'_i\subset N_i$, $N_i$ maps to the unit disk, $N'_i$ maps
to the smaller disk, and the $N'_i$ cover the manifold.
\end{lemma}

\begin{proof}

Choose nested disk charts at each point.  Use second countability of the
separable metric topology to extract a countable subcover of the smaller chart
neighborhoods.
\end{proof}

\begin{lemma}[Initial PL neighborhood]
\label{moise:lem:initial_pl_neighborhood}
\lean{LeanEval.Topology.ClassificationOfSurfaces.initial_pl_neighborhood}

\uses{moise:def:pl_complex_in_space,moise:lem:chart_pair_exhaustion}
For the first chart pair of a two-manifold, there is a finite PL complex
$\mathcal X_1$ in the manifold whose support is a compact two-manifold with
boundary and whose interior subcomplex contains the smaller chart neighborhood.
\end{lemma}

\begin{proof}
\notready
Choose a polygonal disk inside the chart whose interior contains the smaller
disk, triangulate it, and transport the resulting PL complex to the manifold by
the chart homeomorphism.
\end{proof}

\begin{lemma}[Extend a PL complex across one chart]
\label{moise:lem:extend_pl_complex_across_chart}
\lean{LeanEval.Topology.ClassificationOfSurfaces.extend_pl_complex_across_chart}

\uses{moise:def:pl_complex_in_space,moise:def:pl_complex_interior_subcomplex,moise:thm:open_subset_complex,moise:thm:pl_approximation_plane}
Suppose $\mathcal X_n$ is a PL complex in a two-manifold whose support is a
compact two-manifold with boundary and whose interior contains the first $n$
small chart neighborhoods.  Then there is a PL complex $\mathcal X_{n+1}$ with
the same properties for the first $n+1$ chart neighborhoods and with
$\mathcal X_n^\circ\subseteq\mathcal X_{n+1}^\circ$.
\end{lemma}

\begin{proof}
\notready
Use the boundary part of the old support lying in the next chart.  Pull a
neighborhood of this boundary set back to a Euclidean complex, represent the open
subset by a complex using the open-subset theorem, approximate the transition
homeomorphism by a PL homeomorphism using the plane PL approximation theorem,
fit this PL replacement to the old complex, add a finite triangulated disk in the
chart to cover the new smaller chart neighborhood, and subdivide to make all
pieces compatible.
\end{proof}

\begin{theorem}[Moise--Rado triangulation theorem]
\label{moise:thm:rado_triangulation_moise_manifold}
\lean{LeanEval.Topology.ClassificationOfSurfaces.rado_triangulation_moise_two_manifold}

\uses{moise:def:moise_two_manifold,moise:def:pl_complex_in_space,moise:lem:chart_pair_exhaustion,moise:lem:initial_pl_neighborhood,moise:lem:extend_pl_complex_across_chart}
Every Moise two-manifold admits a locally finite PL complex whose support is the
whole manifold.
\end{theorem}

\begin{proof}
\notready

Construct an increasing sequence of PL complexes whose interior subcomplexes
contain progressively more of the chart-pair cover.  The union of the increasing
interior subcomplexes is a locally finite PL complex covering the whole manifold.
\end{proof}

\begin{theorem}[Compact Moise surface is finitely triangulable]
\label{moise:thm:compact_moise_surface_finitely_triangulable}
\lean{LeanEval.Topology.ClassificationOfSurfaces.compact_moise_surface_finitely_triangulable}
\leanok

\uses{moise:thm:rado_triangulation_moise_manifold,moise:lem:compact_locally_finite_complex_finite}
Every compact Moise two-manifold has a finite triangulation.
\end{theorem}

\begin{proof}
\leanok

Apply the Moise--Rado locally finite triangulation theorem.  The support is
compact, so local finiteness implies that only finitely many simplexes occur.
\end{proof}

\section{Boundary surfaces in the Moise route}

\begin{theorem}[Bordered PL approximation]
\label{moise:thm:bordered_pl_approximation}
\lean{LeanEval.Topology.ClassificationOfSurfaces.bordered_pl_approximation}
\leanok

\uses{moise:thm:pl_approximation_plane,moise:def:combinatorial_surface}
The PL approximation theorem remains valid for local half-disk charts and
combinatorial surfaces with boundary, with approximating PL homeomorphisms
respecting the boundary half-plane.
\end{theorem}

\begin{proof}
\leanok
Adapt the plane proof by working in half-plane charts near boundary points and
requiring the one-skeleton approximation to carry boundary edges into the
boundary line.  Interior faces use the ordinary PL Schoenflies argument; boundary
faces use the relative version for half-disks.
\end{proof}

\begin{theorem}[Moise--Rado triangulation for bordered surfaces]
\label{moise:thm:rado_bordered_surface_triangulation}
\lean{LeanEval.Topology.ClassificationOfSurfaces.rado_bordered_surface_triangulation}

\uses{moise:def:moise_two_manifold,moise:thm:rado_triangulation_moise_manifold,moise:thm:bordered_pl_approximation}
Every compact Moise two-manifold with boundary has a finite triangulation whose
boundary is a subcomplex.
\end{theorem}

\begin{proof}
\notready
Either rerun the Rado induction using half-disk chart pairs and the bordered PL
approximation theorem, or triangulate the double and restrict to one side after a
relative subdivision making the boundary a subcomplex.  The direct half-disk
version is the preferred statement because it produces a boundary subcomplex
without an additional symmetry theorem.
\end{proof}

\begin{lemma}[Eval surface to Moise bordered surface]
\label{moise:lem:eval_to_moise_bordered_surface}
\lean{LeanEval.Topology.ClassificationOfSurfaces.eval_surface_to_moise_bordered_surface}

\uses{moise:def:eval_surface,moise:def:moise_two_manifold}
Every compact Eval surface is a compact Moise two-manifold with boundary.
\end{lemma}

\begin{proof}
\notready
Use the half-space manifold charts to obtain local disk or half-disk models.
Derive metrizability and second countability as needed from compactness plus the
finite or countable chart structure available in the mathlib manifold framework.
This is an interface lemma between mathlib's manifold API and Moise's classical
hypotheses.
\end{proof}

\begin{definition}[Finite surface triangulation]
\label{moise:def:finite_surface_triangulation}
\lean{LeanEval.Topology.ClassificationOfSurfaces.FiniteSurfaceTriangulation}

\uses{moise:def:eval_surface,moise:def:euclidean_complex}
A finite surface triangulation of a topological space $S$ is a finite
combinatorial two-manifold with boundary together with a homeomorphism from its
support to $S$.
\end{definition}

\begin{theorem}[Compact Eval surface is finitely triangulable]
\label{moise:thm:compact_eval_surface_finitely_triangulable}
\lean{LeanEval.Topology.ClassificationOfSurfaces.compact_eval_surface_finitely_triangulable}

\uses{moise:def:finite_surface_triangulation,moise:lem:eval_to_moise_bordered_surface,moise:thm:rado_bordered_surface_triangulation}
Every compact connected Eval surface has a finite triangulation with boundary a
subcomplex.
\end{theorem}

\begin{proof}
\notready

Convert the Eval hypotheses to Moise's bordered-surface hypotheses and apply the
compact bordered Rado triangulation theorem.
\end{proof}

\section{From triangulations to cell complexes}

\begin{definition}[Triangulation as a cell complex]
\label{moise:def:triangulation_cell_complex}
\lean{LeanEval.Topology.ClassificationOfSurfaces.FiniteSurfaceTriangulation.toCellComplex}

\uses{moise:def:finite_surface_triangulation,moise:def:surface_cell_complex}
Given a finite surface triangulation, define a surface cell complex with one face
for each triangle, one unoriented edge for each edge of the triangulation, paired
or unpaired oriented darts according to whether the edge is interior or boundary,
and boundary words of length three.
\end{definition}

\begin{lemma}[Triangulation realization matches cell-complex realization]
\label{moise:lem:triangulation_to_cell_complex_realization}
\lean{LeanEval.Topology.ClassificationOfSurfaces.FiniteSurfaceTriangulation.toCellComplex_realization_homeomorphic}

\uses{moise:def:triangulation_cell_complex,moise:def:cell_realization}
The support of a finite surface triangulation is homeomorphic to the realization
of the associated surface cell complex.
\end{lemma}

\begin{proof}
\notready
Each triangle of the triangulation is identified with the corresponding standard
polygonal face.  The simplicial gluing relation on triangle edges agrees with the
cell-complex gluing relation on the associated darts, hence the two quotient
realizations are homeomorphic.
\end{proof}

\begin{theorem}[Finite triangulation to finite cell complex]
\label{moise:thm:finite_triangulation_to_cell_complex}
\lean{LeanEval.Topology.ClassificationOfSurfaces.finite_triangulation_to_cell_complex}

\uses{moise:def:finite_surface_triangulation,moise:def:triangulation_cell_complex,moise:lem:triangulation_to_cell_complex_realization}
If $S$ has a finite surface triangulation, then $S$ is homeomorphic to the
realization of a finite surface cell complex.
\end{theorem}

\begin{proof}

Take the associated cell complex and compose the triangulation homeomorphism with
the homeomorphism between triangulation support and cell-complex realization.
\end{proof}

\section{Quotient realizations and quotient homeomorphisms}

\begin{definition}[Polygonal pre-realization]
\label{moise:def:polygonal_prerealization}
\lean{LeanEval.Topology.ClassificationOfSurfaces.SurfaceCellComplex.PreRealization}

\uses{moise:def:surface_cell_complex}
The pre-realization of $K$ is the finite disjoint union of closed standard
polygons, one for each face of $K$, with marked oriented boundary arcs
corresponding to the cyclic boundary word of that face.
\end{definition}

\begin{definition}[Gluing relation]
\label{moise:def:gluing_relation}
\lean{LeanEval.Topology.ClassificationOfSurfaces.SurfaceCellComplex.GluingRel}

\uses{moise:def:polygonal_prerealization}
The gluing relation on the pre-realization is the least equivalence relation
which identifies paired boundary arcs by orientation-reversing or
orientation-preserving affine parametrizations according to the dart-pairing data
of $K$.
\end{definition}

\begin{theorem}[Mathlib quotient congruence]
\label{moise:thm:mathlib_quotient_congr}
\lean{Homeomorph.Quotient.congr, Homeomorph.Quot.congrRight}
\mathlibok
If a homeomorphism between two spaces carries one quotient relation exactly to
another, it descends to a homeomorphism between the quotient spaces.  If the
underlying space is fixed and two quotient relations are equivalent, the quotient
spaces are homeomorphic.
\end{theorem}

\begin{lemma}[Induced homeomorphism of cell realizations]
\label{moise:lem:cell_realization_quotient_congr}
\lean{LeanEval.Topology.ClassificationOfSurfaces.SurfaceCellComplex.realizationCongr}

\uses{moise:def:cell_realization,moise:def:gluing_relation,moise:thm:mathlib_quotient_congr}
A homeomorphism between polygonal pre-realizations which sends the gluing
relation of $K$ to the gluing relation of $L$ induces a homeomorphism
$|K|\cong |L|$.
\end{lemma}

\begin{proof}

Apply the quotient congruence theorem from mathlib to the homeomorphism of
pre-realizations and the proof that it respects exactly the two gluing relations.
\end{proof}

\begin{definition}[Elementary cell-complex moves]
\label{moise:def:gx_elementary_moves}
\lean{LeanEval.Topology.ClassificationOfSurfaces.SurfaceCellComplex.ElementarySubdivision, LeanEval.Topology.ClassificationOfSurfaces.SurfaceCellComplex.Equivalent}

\uses{moise:def:surface_cell_complex}
The elementary moves are the Gallier--Xu operations: subdivide an edge by
replacing $a,a^{-1}$ with $bc,c^{-1}b^{-1}$, and split a face by inserting a new
diagonal edge.  Equivalence is the equivalence relation generated by these moves.
\end{definition}

\begin{lemma}[Elementary moves preserve realization]
\label{moise:lem:elementary_moves_preserve_realization}
\lean{LeanEval.Topology.ClassificationOfSurfaces.SurfaceCellComplex.elementarySubdivision_realization_homeomorphic}

\uses{moise:def:gx_elementary_moves,moise:lem:cell_realization_quotient_congr}
If $L$ is obtained from $K$ by one elementary Gallier--Xu subdivision move, then
$|K|$ and $|L|$ are homeomorphic.
\end{lemma}

\begin{proof}

For edge subdivision, use the identity on the underlying glued polygonal space
with a reparametrization of the subdivided boundary arc.  For face subdivision,
use the homeomorphism between one polygon and two polygons glued along the new
diagonal.  In both cases, verify the gluing relations and descend by quotient
congruence.
\end{proof}

\begin{theorem}[Equivalent cell complexes have homeomorphic realizations]
\label{moise:thm:equivalent_cell_complexes_homeomorphic}
\lean{LeanEval.Topology.ClassificationOfSurfaces.SurfaceCellComplex.equivalent_realization_homeomorphic}

\uses{moise:def:gx_elementary_moves,moise:lem:elementary_moves_preserve_realization}
If finite surface cell complexes $K$ and $L$ are equivalent in the Gallier--Xu
sense, then $|K|\cong |L|$.
\end{theorem}

\begin{proof}

Induct over the reflexive, symmetric, transitive closure generating equivalence,
composing the homeomorphisms supplied by elementary moves and their inverses.
\end{proof}

\section{Gallier--Xu normal-form reduction}

\begin{definition}[Orientability, contours, and Euler characteristic]
\label{moise:def:gx_invariants}
\lean{LeanEval.Topology.ClassificationOfSurfaces.SurfaceCellComplex.Orientable, LeanEval.Topology.ClassificationOfSurfaces.SurfaceCellComplex.Contours, LeanEval.Topology.ClassificationOfSurfaces.SurfaceCellComplex.eulerChar}

\uses{moise:def:surface_cell_complex}
For a finite surface cell complex, orientability is coherent orientation of
faces, contours are cyclic boundary components, and the Euler characteristic is
$n_0-n_1+n_2$ using the vertices, edges, and faces of the cell complex.
\end{definition}

\begin{lemma}[Elementary moves preserve invariants]
\label{moise:lem:gx_moves_preserve_invariants}
\lean{LeanEval.Topology.ClassificationOfSurfaces.SurfaceCellComplex.elementary_moves_preserve_invariants}

\uses{moise:def:gx_elementary_moves,moise:def:gx_invariants}
The Gallier--Xu elementary moves preserve orientability, the number of contours,
and Euler characteristic.
\end{lemma}

\begin{proof}

Check the two moves.  Edge subdivision increases both the number of vertices and
edges by one.  Face subdivision increases both the number of edges and faces by
one.  Coherent orientations and boundary contours are transported through the
local rewrite.
\end{proof}

\begin{definition}[Canonical cell complexes]
\label{moise:def:canonical_cell_complex}
\lean{LeanEval.Topology.ClassificationOfSurfaces.SurfaceCellComplex.CanonicalOrientable, LeanEval.Topology.ClassificationOfSurfaces.SurfaceCellComplex.CanonicalNonorientable}

\uses{moise:def:surface_cell_complex}
A canonical orientable complex has one face with word
$a_1b_1a_1^{-1}b_1^{-1}\cdots a_pb_pa_p^{-1}b_p^{-1}c_1h_1c_1^{-1}\cdots c_qh_qc_q^{-1}$.
A canonical nonorientable complex has one face with word
$a_1a_1\cdots a_pa_pc_1h_1c_1^{-1}\cdots c_qh_qc_q^{-1}$, with $p\geq 1$.
\end{definition}

\begin{lemma}[Eliminate $aa^{-1}$ strings]
\label{moise:lem:gx_eliminate_inverse_pairs}
\lean{LeanEval.Topology.ClassificationOfSurfaces.SurfaceCellComplex.eliminate_inverse_pairs}

\uses{moise:def:gx_elementary_moves}
Every occurrence of a consecutive string $aa^{-1}$ in a boundary word can be
removed by Gallier--Xu elementary moves.
\end{lemma}

\begin{proof}

Split the face near $aa^{-1}$ using a new edge, contract the resulting length-two
face by inverse edge subdivision, and then remove the redundant face by inverse
face subdivision.
\end{proof}

\begin{lemma}[Vertex reduction]
\label{moise:lem:gx_vertex_reduction}
\lean{LeanEval.Topology.ClassificationOfSurfaces.SurfaceCellComplex.vertex_reduction}

\uses{moise:lem:gx_eliminate_inverse_pairs}
Every connected cell complex is equivalent to one with either a null inner vertex
or exactly one inner vertex, and whose boundary vertices are loops of the form
$(h,c,h^{-1})$.
\end{lemma}

\begin{proof}
\notready

Follow Gallier--Xu Lemma 6.1, Step 2.  Repeatedly shrink an inner vertex by
eliminating an adjacent edge until only one inner vertex remains; then reduce each
boundary vertex to a loop without disturbing already reduced loops.
\end{proof}

\begin{lemma}[Reduce to one face and introduce crosscaps]
\label{moise:lem:gx_one_face_crosscaps}
\lean{LeanEval.Topology.ClassificationOfSurfaces.SurfaceCellComplex.one_face_and_crosscaps}

\uses{moise:lem:gx_vertex_reduction}
After vertex reduction, every connected cell complex is equivalent to one with a
single face, and repeated occurrences $aXaY$ may be replaced by adjacent
crosscaps.
\end{lemma}

\begin{proof}
\notready

Use inverse face subdivision to merge adjacent faces sharing an edge until one
face remains.  Then use the Gallier--Xu pseudo-rewrite
$aXaY\simeq bbY^{-1}X$ to introduce crosscaps while preserving loops and the
reduced vertex structure.
\end{proof}

\begin{lemma}[Introduce handles]
\label{moise:lem:gx_introduce_handles}
\lean{LeanEval.Topology.ClassificationOfSurfaces.SurfaceCellComplex.introduce_handles}

\uses{moise:lem:gx_one_face_crosscaps}
Boundary words containing interlaced inverse pairs $aUbVa^{-1}Xb^{-1}Y$ are
equivalent to words containing a handle $cdc^{-1}d^{-1}$.
\end{lemma}

\begin{proof}
\notready

Formalize the pseudo-rewrite
$aUVa^{-1}X\simeq bVUb^{-1}X$ and apply it twice to transform the interlaced pair
into a handle block.
\end{proof}

\begin{lemma}[Handle plus crosscap becomes three crosscaps]
\label{moise:lem:gx_handle_crosscap_to_crosscaps}
\lean{LeanEval.Topology.ClassificationOfSurfaces.SurfaceCellComplex.handle_crosscap_to_three_crosscaps}

\uses{moise:lem:gx_one_face_crosscaps,moise:lem:gx_introduce_handles}
A word containing one crosscap and one handle is equivalent to a word containing
three crosscaps instead.
\end{lemma}

\begin{proof}
\notready

Use the Gallier--Xu pseudo-rewrite $aaXY\simeq bYbX^{-1}$ repeatedly on the word
$aaXbcb^{-1}c^{-1}Y$ to obtain three adjacent crosscap blocks.
\end{proof}

\begin{lemma}[Group boundary loops]
\label{moise:lem:gx_group_loops}
\lean{LeanEval.Topology.ClassificationOfSurfaces.SurfaceCellComplex.group_boundary_loops}

\uses{moise:lem:gx_introduce_handles}
All loop blocks $chc^{-1}$ can be moved to the end of the single boundary word
without changing the handle or crosscap blocks.
\end{lemma}

\begin{proof}
\notready

Use the pseudo-rewrite $aUVa^{-1}X\simeq bVUb^{-1}X$ to commute any two loop
blocks until all loops are grouped together.
\end{proof}

\begin{theorem}[Every cell complex has canonical normal form]
\label{moise:thm:gx_normal_form}
\lean{LeanEval.Topology.ClassificationOfSurfaces.SurfaceCellComplex.equivalent_canonical}

\uses{moise:def:canonical_cell_complex,moise:lem:gx_vertex_reduction,moise:lem:gx_one_face_crosscaps,moise:lem:gx_introduce_handles,moise:lem:gx_handle_crosscap_to_crosscaps,moise:lem:gx_group_loops}
Every finite connected surface cell complex is equivalent to a canonical
orientable or canonical nonorientable cell complex.
\end{theorem}

\begin{proof}

Apply the Gallier--Xu sequence of reductions: remove inverse pairs, reduce
vertices, merge faces, introduce crosscaps, introduce handles, convert mixed
handle-crosscap words to all-crosscap form, and group boundary loops.
\end{proof}

\section{Normal-form quotient representatives and final assembly}

\begin{definition}[Eval quotient representatives]
\label{moise:def:eval_quotient_representatives}
\lean{LeanEval.Topology.ClassificationOfSurfaces.OrientableRel, LeanEval.Topology.ClassificationOfSurfaces.NonOrientableRel}
The Eval representatives are quotient spaces determined by the orientable and
nonorientable normal-form polygon boundary relations.  They are the exact target
spaces named in the Lean Eval classification statement.
\end{definition}

\begin{lemma}[Canonical realization equals Eval quotient]
\label{moise:lem:canonical_realization_eval_quotient}
\lean{LeanEval.Topology.ClassificationOfSurfaces.canonical_realization_homeomorphic_eval_quotient}

\uses{moise:def:canonical_cell_complex,moise:def:eval_quotient_representatives,moise:lem:cell_realization_quotient_congr}
The realization of a canonical orientable or nonorientable cell complex is
homeomorphic to the corresponding Eval quotient representative.
\end{lemma}

\begin{proof}

Identify the one polygon used by the canonical cell complex with the polygon used
in the Eval quotient definition and prove that the two generated boundary
relations coincide.  Apply quotient congruence.
\end{proof}

\begin{theorem}[Cell-complex classification into Eval representatives]
\label{moise:thm:cell_complex_has_eval_representative}
\lean{LeanEval.Topology.ClassificationOfSurfaces.cell_complex_has_eval_representative}

\uses{moise:thm:equivalent_cell_complexes_homeomorphic,moise:thm:gx_normal_form,moise:lem:canonical_realization_eval_quotient}
For every finite connected surface cell complex $K$, the realization $|K|$ is
homeomorphic to the sphere representative, or to an orientable Eval quotient
representative, or to a nonorientable Eval quotient representative.
\end{theorem}

\begin{proof}

Choose a canonical cell complex equivalent to $K$.  Equivalent complexes have
homeomorphic realizations, and canonical realizations are the Eval quotient
representatives.
\end{proof}

\begin{theorem}[Topological classification of compact Eval surfaces]
\label{moise:thm:eval_classification_of_surfaces}
\lean{LeanEval.Topology.ClassificationOfSurfaces.topological_classification_of_surfaces}

\uses{moise:thm:compact_surface_cellulable,moise:thm:cell_complex_has_eval_representative}
Every compact connected Hausdorff topological $2$-manifold with boundary in the
Lean Eval sense is homeomorphic to one of the normal-form quotient spaces listed
in the Eval problem.
\end{theorem}

\begin{proof}

Apply the topological bridge to obtain $S\cong |K|$.  Apply the cell-complex
classification to $K$.  Compose the resulting homeomorphisms.
\end{proof}
